%% toc: P2 22 22.8
\section{Where \(L_{\omega_1\omega}\) is developed further, and how it differs from \(L_{\infty\omega}\)}
\label{sec:p02-w02-p02-039-lomega-treatments}

\Cref{sec:p02-infinitary-generalized} fixes the grammar and satisfaction
clauses of \(L_{\omega_1\omega}\) in \cref{def:p02-lomega-one-omega}, states
its basic calibration in \cref{thm:p02-lomega-one-omega-calibration}, and
records the infinitary proof boundary in \cref{thm:p02-karp-completeness}.
There the calibration theorem and the Scott construction that follows the
Karp boundary are supported by citation only.  This section does two things.
First, it locates claim by claim the places in the handbook where those
statements are proved in full or developed further
(\cref{tab:p02-w02-p02-039-treatments}).  Second, it separates
\(L_{\omega_1\omega}\) from the unrestricted infinitary language
\(L_{\infty\omega}\): \cref{thm:p02-w02-p02-039-proper-fragment} proves that,
already in a signature with one binary relation symbol, some
\(L_{\infty\omega}\)-sentence has a model class that no
\(L_{\omega_1\omega}\)-sentence has.

\emph{Inputs.}  \Cref{def:p02-lomega-one-omega}; the finite partial
isomorphisms, potential isomorphisms and the language \(L_{\infty\omega}\) of
\cref{def:p08-bvb501-potential-isomorphism}; and the characteristic sentences
of \cref{thm:p08-bvb501-potential-isomorphism}.  \emph{Conclusions.}  A map
from each claim of \cref{sec:p02-infinitary-generalized} to its fuller
treatment, and a proof that \(L_{\omega_1\omega}\) is a proper fragment of
\(L_{\infty\omega}\) in expressive power.  \emph{Ambient assumptions.}
Metatheory is ZFC set theory.  Choice is used once, in
\cref{thm:p02-w02-p02-039-formula-set}, to know that a countable set of
countable ordinals has a countable supremum.  Structures have nonempty
domains, as in \cref{def:p02-lomega-one-omega}.

\begin{hbconvention}{def:p02-w02-p02-039-notation}{Notation across the infinitary treatments}
\begin{enumerate}
\item The language written \(L_{\omega_1\omega}\) in
      \cref{def:p02-lomega-one-omega} and in
      \cref{def:p08-bvb501-potential-isomorphism} is the language written
      \(L_{\omega_1,\omega}\) in \cref{def:p03-morley-infinitary}.  All
      three admit conjunctions and disjunctions of countable sets of
      formulas whose free variables form a finite set in total, and only
      finite strings of individual quantifiers.  The comma is a
      typographical difference only.
\item \(L_{\infty\omega}\) is the language of
      \cref{def:p08-bvb501-potential-isomorphism}: conjunctions and
      disjunctions of arbitrary sets of formulas, again with a finite total
      free-variable set and finite quantifier strings.  Its satisfaction
      clauses are those of \cref{def:p02-lomega-one-omega} with the word
      ``countable'' removed.  Every \(L_{\omega_1\omega}\)-formula is an
      \(L_{\infty\omega}\)-formula and is satisfied by the same structures
      and assignments.
\item Two sentences of the same signature are \emph{equivalent} when they
      have exactly the same models.
\item In this section \(\delta_M\) is always the characteristic sentence of
      \cref{thm:p08-bvb501-potential-isomorphism}: an
      \(L_{\infty\omega}\)-sentence whose models are exactly the structures
      potentially isomorphic to \(M\).  The same letter in
      \cref{thm:p03-morley-scott-stabilization} denotes a countable ordinal,
      and the countable-model sentence there is written \(\sigma_M\).
\item For an ordinal \(\alpha\geq1\), \((\alpha,<)\) is the structure
      whose domain is the set of ordinals below \(\alpha\) and whose single
      binary relation symbol \(<\) is interpreted by the ordinal order.
\end{enumerate}
\end{hbconvention}

\subsection{Claim-by-claim location of the fuller treatments}
\label{sec:p02-w02-p02-039-locations}

\begingroup\small
\begin{longtable}{@{}p{.28\linewidth}p{.38\linewidth}p{.22\linewidth}@{}}
\caption{Claims about \(L_{\omega_1\omega}\) in
\cref{sec:p02-infinitary-generalized} and where each is developed further.}
\label{tab:p02-w02-p02-039-treatments}\\
\toprule
Claim & Fuller treatment & Standing there \\ \midrule
\endfirsthead
\toprule
Claim & Fuller treatment & Standing there \\ \midrule
\endhead
Grammar with the finite total free-variable condition, and the satisfaction
clauses (\cref{def:p02-lomega-one-omega}). &
\Cref{def:p03-morley-infinitary} adds well-founded formula trees and a
countable-ordinal quantifier rank; \cref{def:p08-bvb501-potential-isomorphism}
places the language inside \(L_{\infty\omega}\). &
Definitions. \\
\addlinespace
Failure of compactness (first sentence of
\cref{thm:p02-lomega-one-omega-calibration}). &
The counterexample displayed after that theorem; its verification is
written out below this table. &
Proved in this chapter. \\
\addlinespace
For a countable structure in a countable signature, a sentence whose
countable models are exactly its isomorphic copies
(\cref{thm:p02-lomega-one-omega-calibration}, and the construction
paragraph after \cref{thm:p02-karp-completeness}). &
\Cref{thm:p03-morley-scott-sentence}, built from
\cref{def:p03-morley-approximants,def:p03-morley-scott-formulas} and
\cref{thm:p03-morley-scott-stabilization}; independently, the last
clause of \cref{thm:p08-bvb501-potential-isomorphism}. &
Proved in the annex for each of the two routes. \\
\addlinespace
The restriction of that characterization to countable models. &
The final sentence of \cref{thm:p03-morley-scott-sentence}; the
\((\mathbb Q,<)\) and \((\mathbb R,<)\) comparison after
\cref{thm:p08-bvb501-first-order-game}, which exhibits potentially
isomorphic structures of different cardinalities. &
Proved (the comparison is verified in the proof of
\cref{thm:p08-bvb501-potential-isomorphism}). \\
\addlinespace
Agreement of approximant-equivalent tuples on ranked formulas, and
isomorphism of countable structures as agreement at every countable stage. &
\Cref{thm:p03-morley-approximant-properties,thm:p03-morley-approximant-isomorphism}. &
Proved in the annex. \\
\addlinespace
Truth of a fixed formula as a function of a coded countable structure. &
\Cref{thm:p03-morley-borel-truth}. &
Proved in the annex. \\
\addlinespace
Agreement on complete \(L_{\infty\omega}\)-theories as potential
isomorphism. &
\Cref{thm:p08-bvb501-potential-isomorphism}. &
Proved in the annex. \\
\addlinespace
\(L_{\omega_1\omega}\) is not the unrestricted language. &
\Cref{thm:p02-w02-p02-039-proper-fragment} below. &
Proved in the annex. \\
\addlinespace
Karp's completeness boundary for well-founded infinitary derivations
(\cref{def:p02-infinitary-derivation,thm:p02-karp-completeness}). &
No fuller treatment exists in this handbook.  The attribution of the
completeness theorem for a Gentzen-type infinitary calculus to Karp and to
Lopez-Escobar is recorded in \cite[ch.~A.7, \S1, p.~236]{Makkai1977}. &
Cited only. \\
\addlinespace
Role in abstract model theory. &
\Cref{sec:p03-abstract-landmarks}, the hypotheses of
\cref{thm:p03-lindstrom}, and the
\hyperref[entry:p14-directory-infinitary]{directory entry for infinitary and
generalized-quantifier languages}. &
Comparison; Lindstr\"om's theorem is proved in the annex. \\
\bottomrule
\end{longtable}
\endgroup

The compactness counterexample of \cref{sec:p02-infinitary-generalized} is
verified as follows.  Let \(\Gamma\) be the set
\(\{\bigvee_{n<\omega}P_n\}\cup\{\neg P_n:n<\omega\}\) of
\(L_{\omega_1\omega}\)-sentences over the nullary relation symbols \(P_n\).
A finite \(\Gamma_0\subseteq\Gamma\) contains \(\neg P_n\) only for \(n\) in
some finite set \(N_0\subseteq\omega\).  Choose \(m\in\omega\setminus N_0\)
and let \(\mathcal M\) make \(P_m\) true and every other \(P_n\) false.
Then \(\mathcal M\) satisfies every \(\neg P_n\) with \(n\in N_0\), and it
satisfies the disjunction by the disjunction clause of
\cref{def:p02-lomega-one-omega} with witness \(P_m\).  A model of all of
\(\Gamma\) would make every \(P_n\) false and, by the same clause, some
\(P_n\) true.  The set \(\Gamma\) is countable, so even countable
compactness fails.

\subsection{\texorpdfstring{\(L_{\omega_1\omega}\)}{L omega1 omega} is a proper fragment of \texorpdfstring{\(L_{\infty\omega}\)}{L infinity omega}}
\label{sec:p02-w02-p02-039-proper-fragment}

The positive results in \cref{tab:p02-w02-p02-039-treatments} depend on two
countability conditions at once: the connective index sets are countable,
and the structures being characterized are countable.  The unrestricted
language drops the first condition.  The next three results show that this
is a genuine gain in expressive power and not a change of presentation.

\begin{hblemma}{thm:p02-w02-p02-039-formula-set}{proved=annex:p02-w02-p02-039-formula-set}{The formulas of \(L_{\omega_1\omega}\) over a fixed signature form a set}
Let \(\tau\) be a signature that is a set, with a countable set of
variables.  Regard each formula as a set built from its immediate
constituents by a fixed set-theoretic operation, so that \(\bigwedge\Phi\)
and \(\bigvee\Phi\) are built from the set \(\Phi\).  Then the class of
\(L_{\omega_1\omega}\)-formulas over \(\tau\) is a set, and so is the class
of its sentences.
\end{hblemma}

\begin{hblemma}{thm:p02-w02-p02-039-ordinal-rigidity}{proved=annex:p02-w02-p02-039-rigid}{Ordinals are rigid under potential isomorphism}
Let \(\alpha,\beta\geq1\) be ordinals and let \(I\) be a potential
isomorphism between \((\alpha,<)\) and \((\beta,<)\) in the sense of
\cref{def:p08-bvb501-potential-isomorphism}.  Then every map in \(I\) is
the identity on its domain, and \(\alpha=\beta\).
\end{hblemma}

\begin{hbtheorem}{thm:p02-w02-p02-039-proper-fragment}{proved=annex:p02-w02-p02-039-proper-fragment}{\(L_{\omega_1\omega}\) is strictly weaker than \(L_{\infty\omega}\)}
Work in the signature \(\{<\}\) with one binary relation symbol.
\begin{enumerate}
\item For every countable ordinal \(\alpha\geq1\), the sentence
      \(\delta_{(\alpha,<)}\) is equivalent to an
      \(L_{\omega_1\omega}\)-sentence.
\item There is an ordinal \(\xi\) such that for every ordinal
      \(\alpha\geq\xi\), the \(L_{\infty\omega}\)-sentence
      \(\delta_{(\alpha,<)}\) is equivalent to no
      \(L_{\omega_1\omega}\)-sentence.  Every such \(\alpha\) is
      uncountable.
\end{enumerate}
Consequently \(L_{\omega_1\omega}\) and \(L_{\infty\omega}\) define
different families of model classes over \(\{<\}\), and the
\(L_{\infty\omega}\)-sentences over \(\{<\}\) do not form a set.
\end{hbtheorem}

The proof is a counting argument: the sentences \(\delta_{(\alpha,<)}\) for
distinct ordinals \(\alpha\) have distinct model classes by
\cref{thm:p02-w02-p02-039-ordinal-rigidity}, so there is a proper class of
them, while \cref{thm:p02-w02-p02-039-formula-set} leaves only a set of
\(L_{\omega_1\omega}\)-sentences to match them.  The same set-versus-class
contrast for complete theories is noted after
\cref{thm:p08-bvb501-first-order-game}.  The argument does not name a
specific \(\alpha\); it bounds the ordinals that can be matched.

\begin{hbexample}{def:p02-w02-p02-039-example}{Why no potential isomorphism links \((\omega,<)\) and \((\omega+1,<)\)}
The map \(f=\{0\mapsto1\}\) from \((\omega,<)\) to \((\omega+1,<)\) is a
finite partial isomorphism: a one-point map preserves and reflects the
atomic formulas \(x=x\) and \(x<x\).  Suppose \(f\) belonged to a potential
isomorphism \(I\).  The element \(0\) of \(\omega+1\) would have to enter
the range of some extension \(g\in I\) of \(f\), say \(g(c)=0\).  Since
\(g(c)=0<1=g(0)\) and \(g\) reflects \(<\), we get \(c<0\) in \(\omega\),
which is impossible.  This is the step that drives
\cref{thm:p02-w02-p02-039-ordinal-rigidity}: a map that moves some point
upward fails the back condition, and one that moves a point downward fails
the forth condition.  The two structures are also separated by a
first-order sentence, \(\exists x\,\forall y\,(y<x\lor y=x)\), which holds
only in \((\omega+1,<)\); the example illustrates the mechanism, not a gap
between the languages.  For countable ordinals no gap exists at the level of
single characteristic sentences, by clause~1 of
\cref{thm:p02-w02-p02-039-proper-fragment}.
\end{hbexample}

\begin{hbwarning}{def:p02-w02-p02-039-not-unrestricted}{Do not identify \(L_{\omega_1\omega}\) with an unrestricted infinitary language}
\begin{enumerate}
\item \(L_{\omega_1\omega}\) admits only countable conjunctions and
      disjunctions.  By \cref{thm:p02-w02-p02-039-proper-fragment} the
      unrestricted language \(L_{\infty\omega}\) defines model classes that
      \(L_{\omega_1\omega}\) does not, so a result proved for one is not
      thereby a result for the other.
\item The unrestricted language is still limited.  Potentially isomorphic
      structures satisfy the same \(L_{\infty\omega}\)-sentences
      (\cref{thm:p08-bvb501-potential-isomorphism}), and \((\mathbb Q,<)\)
      and \((\mathbb R,<)\) are potentially isomorphic.  Hence no
      \(L_{\infty\omega}\)-sentence, and in particular no
      \(L_{\omega_1\omega}\)-sentence, is true in one of them and false in
      the other, although they are not isomorphic.
\item Both languages keep quantifier strings finite; that is the meaning of
      the second subscript \(\omega\) in \cref{def:p02-lomega-one-omega}.
      Nothing in this section concerns languages that quantify over
      infinitely many variables at once.
\item The Karp boundary of \cref{thm:p02-karp-completeness} concerns the
      Karp calculus for \(L_{\omega_1\omega}\) with well-founded infinitary
      derivations (\cref{def:p02-infinitary-derivation}).  It is not
      extended here to \(L_{\infty\omega}\), and it does not restore
      compactness for either language.
\item The Scott sentences of \cref{thm:p03-morley-scott-sentence} and the
      countable clause of \cref{thm:p08-bvb501-potential-isomorphism}
      characterize a countable structure among countable structures only.
      Clause~2 of \cref{thm:p02-w02-p02-039-proper-fragment} shows that
      for all sufficiently large ordinals \(\alpha\) no
      \(L_{\omega_1\omega}\)-sentence has as its models exactly the
      structures potentially isomorphic to \((\alpha,<)\), although
      \(L_{\infty\omega}\) has such a sentence.
\end{enumerate}
\end{hbwarning}

\subsection{Exact sources}
\label{sec:p02-w02-p02-039-sources}

The original sources for the claims of \cref{sec:p02-infinitary-generalized}
are \textcite{Karp1964}, for the infinitary proof calculus and its
completeness, and \textcite{Scott1965}, for the model theory of
\(L_{\omega_1\omega}\) and the characterization of countable structures.
The proofs referred to in \cref{tab:p02-w02-p02-039-treatments} and the
three results of \cref{sec:p02-w02-p02-039-proper-fragment} use in-book
definitions only and do not depend on either source.  The following
survey locators record, in a single reference work, each claim as this
section uses it.
\begin{enumerate}
\item \emph{Grammar.}  The formation rule that admits the conjunction and
      the disjunction of any countable set of formulas, and the reading of
      the subscripts as ``fewer than \(\omega_1\) connective arguments,
      fewer than \(\omega\) quantified variables in a string'', are given in
      \cite[ch.~A.1, \S5.4, pp.~43--44]{Barwise1977FirstOrder}.
\item \emph{Compactness and completeness boundary.}  The same passage
      records that \(L_{\omega_1\omega}\) satisfies the L\"owenheim--Skolem
      theorem but not the compactness theorem, and that a completeness
      theorem for countable theories needs an infinitary rule of proof
      \cite[ch.~A.1, \S5.4, p.~44]{Barwise1977FirstOrder}.  The
      compactness failure is verified directly after
      \cref{tab:p02-w02-p02-039-treatments}; the L\"owenheim--Skolem
      statement is reported here and is not used in this section.
\item \emph{Karp's calculus.}  The completeness theorem for
      \(L_{\omega_1\omega}\), in the form that every valid formula has a
      derivation in a Gentzen-type infinitary system, is attributed to
      \textcite{Karp1964} and \textcite{LopezEscobar1965} in
      \cite[ch.~A.7, \S1, p.~236]{Makkai1977}; the same page records
      that a countable admissible fragment strengthens this by choosing
      the derivation inside the admissible set.  Karp's book is also named
      there as the first systematic investigation of infinitary languages,
      and \textcite{Scott1965} as an early account of the model theory of
      \(L_{\omega_1\omega}\) \cite[ch.~A.7, \S1, p.~235]{Makkai1977}.
\item \emph{Scott sentences.}  The isomorphism theorem of
      \textcite{Scott1965} is obtained as a direct consequence of a
      theorem of Vaught: every class of countable structures that is
      \(\Sigma^1_1\) over \(L_{\omega_1\omega}\) is a union of \(\aleph_1\)
      classes, each defined by an \(L_{\omega_1\omega}\)-sentence
      \cite[ch.~A.7, \S8, Theorem~8.11 and the remark after it,
      p.~276]{Makkai1977}.  The in-book proofs listed in
      \cref{tab:p02-w02-p02-039-treatments} follow the approximant route
      of \cref{thm:p03-morley-scott-sentence} and the potential-isomorphism
      route of \cref{thm:p08-bvb501-potential-isomorphism}, not this
      descriptive-set-theoretic route.
\item \emph{A characterizing sentence.}  An \(L_{\omega_1\omega}\)-sentence
      that pins down the standard natural numbers up to isomorphism, the
      conjunction of a diagram with the countable disjunction saying that
      each element is named by a numeral, is displayed in
      \cite[ch.~A.7, \S5, p.~254]{Makkai1977}.
\end{enumerate}
